\documentclass[11pt]{article}

\usepackage[margin=1in]{geometry}
\usepackage{amsmath,amssymb,amsthm}
\usepackage{empheq}
\usepackage{graphicx}
\usepackage{booktabs}
\usepackage{microtype}
\usepackage[hidelinks]{hyperref}
\usepackage{siunitx}
\usepackage{authblk}

\newcommand{\diff}{\mathrm{d}}
\newcommand{\eff}{\mathrm{eff}}
\newcommand{\BP}{\mathrm{BP}}
\newcommand{\EP}{\mathrm{EP}}
\newcommand{\msd}{\mathrm{MSD}}

\title{Chromatin first-passage as a constraint on developmental
positional information, and its consequences for single-molecule readout
in gastruloids}

\author{[Author]}
\affil{[Affiliation]}

\date{\today}

\begin{document}

\maketitle

\begin{abstract}
Lammers, Flamholz, and Garcia~\cite{Lammers2023} recently showed that the
information rate driving cellular decisions is bounded by a three-way
tradeoff between transcriptional sharpness, precision, and specificity, and
that the early \textit{Drosophila} embryo and mammalian cells satisfy this
bound through fundamentally different reliance on energy dissipation,
indexed by the noncognate-to-cognate factor concentration ratio $w/c$. We
argue here that a parallel and physically distinct bound arises from
chromatin polymer first-passage: even when a transcription factor (TF) has
correctly identified its cognate enhancer, productive transcription
requires the enhancer to come into contact with the promoter, and the
distribution of contact times is set by the chromatin mean-square
displacement (MSD) exponent~$\alpha$ and the genomic separation~$s$. We
write down the resulting compound noise budget combining diffusive TF
counting (Berg--Purcell), polymer-FPT noise on enhancer--promoter contact,
and the Lammers proofreading term, and analyze its scaling with~$\alpha$.
The same fly-vs-mouse asymmetry that Lammers identifies on the $w/c$ axis
appears \emph{simultaneously} on the chromatin axis: \textit{Drosophila}
nuclear cycle 14 sits comfortably because chromatin is pre-TAD and
Rouse-like ($\alpha\!\approx\!0.5$); the mouse early embryo sits at the
bound only by replacing fast single-locus first-passage with multi-enhancer
redundancy, paracrine cell--cell averaging, and intra-TAD geometry. We
specialize to 2D mouse gastruloids, where the morphogen gradient is
self-organized rather than imposed, and propose a single-molecule
measurement program based on H2B and Sox2/Bra tracking that decomposes the
readout-noise budget term-by-term and discriminates reaction--diffusion
from localized-source mechanisms of symmetry breaking.
\end{abstract}

\section{Introduction}

The classical Berg--Purcell bound on diffusion-limited concentration sensing
states that a receptor of size $a$ measuring a ligand at free concentration
$c$ for a time $T$ has fractional concentration error
\begin{equation}
  \left(\frac{\delta c}{c}\right)^{\!2}_{\BP}
  \;\gtrsim\; \frac{1}{a\,D\,c\,T}.
  \label{eq:bp}
\end{equation}
Gregor and colleagues~\cite{Gregor2007a,Gregor2007b} applied this bound to
the readout of the Bicoid (Bcd) gradient in the syncytial \textit{Drosophila}
blastoderm and showed that the embryo positions the \textit{hunchback}
expression boundary at $\sim$1\% egg length, essentially saturating
\eqref{eq:bp} for the available cycle-14 integration time of $T \approx
\SI{180}{\second}$. Recent work has extended this picture in two directions
that motivate the present paper.

First, Lammers, Flamholz, and Garcia~\cite{Lammers2023} introduced an
information-rate (IR) framework based on a kinetic, multi-state model of
transcriptional bursting. Their bound on the time required to distinguish
two TF concentrations differing by $\delta c/c$ takes the form
\begin{equation}
  \overline{T} \;\geq\; \frac{\ln\!\left(\frac{1-\varepsilon}{\varepsilon}\right)(1-2\varepsilon)}{\mathrm{IR}},
  \qquad
  \mathrm{IR} \;=\; \tfrac{1}{2}\!\left(\tfrac{\delta c}{c^{*}}\right)^{\!2}\, s^{2}\, p^{2},
  \label{eq:lammers}
\end{equation}
where~$s$ and~$p$ are sharpness and precision of the gene's input--output
function. They show that at equilibrium, sharpness is bounded by the number
of TF binding sites and precision by intrinsic gene-circuit noise; the
specificity~$f$ relating cognate to noncognate binding is bounded by the
intrinsic affinity ratio~$\alpha_{Lammers}$. In the high-interference
regime where the noncognate factor concentration~$w$ is large compared with
$\alpha_{Lammers}\,c$, equilibrium gene circuits cannot drive timely
decisions; energy dissipation through kinetic-proofreading-like cycles is
required to break the sharpness--precision--specificity tradeoff. Lammers
\textit{et al.} estimate $w/c\approx 47$ for the early \textit{Drosophila}
embryo (low interference; equilibrium suffices), $w/c\le 698$ for
\textit{C.~elegans}, and $w/c\le 1426$ for mouse (high interference;
nonequilibrium activator proofreading required).

Second, Br\"uckner and colleagues~\cite{Bruckner2023} measured directly,
using paired live-imaging of enhancer and promoter loci, that
enhancer--promoter (EP) encounter times in fly nuclei show \emph{anomalous
scaling with genomic separation}, weaker than predicted by the standard
fractional-Langevin / Rouse polymer scaling. Hansen-lab two-color tracking
of CTCF-anchored loops in mESCs~\cite{Gabriele2022,Mach2022} measured
mammalian-interphase EP-equivalent contacts as rare and short:
minute-to-tens-of-minute lifetimes occupying only a few percent of the
cell cycle.

These two threads --- Lammers' kinetic bound and the empirical chromatin
dynamics --- bear on the same problem (timely transcriptional decisions
under physical noise) but address physically distinct bottlenecks.
Lammers~\cite{Lammers2023} models the locus as a single point, with
specificity-vs-precision tradeoffs against an interfering noncognate factor
pool; the chromatin polymer in which the locus is embedded does not enter.
Br\"uckner~\cite{Bruckner2023} measures the polymer dynamics but does not
derive a noise budget. Lucas \textit{et al.}~\cite{Lucas2018} explicitly
note, on the basis of MS2 imaging at \textit{hb}, that Berg--Purcell
predictions do not fully account for the observed precision and that
``additional mechanisms beyond simple diffusive readout'' must be at play.

Here we close that gap. We treat chromatin polymer first-passage as a
distinct physical constraint on transcriptional decision time, write down
the resulting compound noise budget, and show that it identifies a second
fly-vs-mouse asymmetry --- on the chromatin MSD exponent~$\alpha$ ---
running parallel to Lammers' $w/c$ asymmetry. The two constraints are
independent: a system can satisfy one and fail the other. Both must be
satisfied for the readout to be physically realizable in the available
window. We then specialize to 2D mouse gastruloids, in which both axes are
experimentally accessible by single-molecule tracking of histone H2B
(chromatin axis) and of Sox2/Bra (TF kinetics axis), and propose a
measurement program decomposing the readout-noise budget term-by-term.

\section{Related work}
\label{sec:related}

\paragraph{Lammers, Flamholz, Garcia 2023~\cite{Lammers2023}.}
The most direct precedent. Frames cellular decisions as an information-rate
problem with a precision/sharpness/specificity tradeoff, and derives the
fly-vs-mouse split from the noncognate-factor interference ratio~$w/c$. We
adopt their decision-time formulation~\eqref{eq:lammers} and treat our
polymer term as an additional, physically independent constraint operating
in parallel. Their model takes the gene locus as a structureless point;
chromatin dynamics enter neither the sharpness, precision, nor specificity
budget. They explicitly note (\textit{op.~cit.}) that ``transcription is
also thought to require multiple molecular steps beyond activator binding
itself, such as the localization of key general transcription factors to
the gene locus,'' which the present framework formalizes.

\paragraph{Br\"uckner \textit{et al.} 2023~\cite{Bruckner2023}.} Provides
the empirical foundation for treating chromatin first-passage as an
independent constraint. Two-color live-imaging of EP-pair separation in
\textit{Drosophila} nuclei shows ``the coexistence of a compact globular
organization and fast subdiffusive dynamics. These combined features cause
an anomalous scaling of polymer relaxation times with genomic separation
leading to long-ranged correlations. Thus, encounter times of DNA loci are
much less dependent on genomic distance than predicted by existing polymer
models, with potential consequences for eukaryotic gene expression.''
This finding pushes effective $\tau_c$ \emph{down} relative to the naive
Rouse scaling, which favors the compatibility of the polymer term with
\textit{Drosophila} nc14 readout.

\paragraph{Shi, Shin, Thirumalai 2025~\cite{Shi2025}.} A theoretical
treatment that resolves the Br\"uckner conundrum and reframes the polymer
side of the present bound. Using the HIPPS-DIMES maximum-entropy
reconstruction, they show that the static three-dimensional contact map
(Hi-C / Micro-C) determines a connectivity matrix~$\mathbf{K}$ whose
eigendecomposition predicts the two-point relaxation times $\tau$, the
mean first-passage times $\tau_c$, and single-locus diffusivities ---
quantitatively reproducing the Br\"uckner $\tau\sim s^{0.7\text{--}0.8}$
scaling and the experimental $\tau_c\sim\langle r\rangle^{3.4}$. Three
findings bear directly on this paper. First, randomly shuffling the WT
contact map (preserving polymer connectivity) destroys the scaling:
the rapid dynamics is a property of the specific cohesin- and TF-organized
contact pattern, not of any homopolymer property. Second, even an
\emph{effective-homopolymer} reduction of the WT distance map gives
$\tau\sim s^{1.1}$ and $\tau\sim r^{4}$, distinct from the WT exponents,
implying that the closed-form fractional-Langevin scaling
\eqref{eq:tau_c} is structurally inadequate for chromatin and should be
read as a homopolymer reference rather than as the operative law. Third,
acute cohesin depletion ($\Delta$RAD21) reorganizes $\mathbf{K}$ in a
locus-specific, signed manner: $\tau_{ij}$ \emph{increases} for intra-TAD
pairs (because $\langle r\rangle$ grows) and \emph{decreases} for
TAD-boundary pairs, while ensemble single-locus diffusivity rises
$\sim$20--40\% with the local fractional-Langevin exponent~$\alpha$
largely preserved. Their derivation of a polymer-FPT Hill law,
$\hat k = 1/(1+(\tau_0/\tau_d)p_c^{-\theta})$, supplies a mechanistic
bridge from~$\tau_c$ to a hypersensitive transcriptional output that is
independent of, and complementary to, the Xiao~\cite{Xiao2021}
futile-cycle bistability.

\paragraph{Kim \& Larson 2025~\cite{KimLarson2025}.} The closest
conceptual precedent for the coupled framework derived in
Section~\ref{sec:gated}. Reviews transcriptional bursting models in
development and explicitly extends the two-state telegraph model with a
kinetic-proofreading step (their Fig.~3, p.~1000), arriving at the
specificity expression
$\psi \approx 1 - E/M$ with $E = k_f k_s/(k_{\mathrm{off}}k_p)$ and $M$ the
ratio of nonspecific to specific binding. Their treatment formalizes the
speed--efficiency--specificity tradeoff that Gregor's timing constraint
implies for development and is, to our knowledge, the only published
extension of the Gregor framework with explicit kinetic proofreading.
Their proofreading model, however, assumes that the substrate (the TF)
is at the locus and that the cycle runs freely once initiated; the
chromatin-polymer step that brings enhancer and promoter into productive
contact, and the finite contact dwell time that gates the proofreading
cycle, do not enter their derivation. Kim \& Larson~\cite{KimLarson2025}
explicitly identify this as an open question (their Question~1, p.~1007:
``How do different chromatin states and 3D genome organization influence
transcription bursting kinetics?''). Section~\ref{sec:gated} of the
present work fills that gap by formally coupling polymer first-passage
to the proofreading cycle and showing that the resulting bound is a
threshold rather than a smooth tradeoff.

\paragraph{Garcia, Upadhyaya, Hager 2021~\cite{Garcia2021NAR}.}
Establishes that transcription-factor residence-time distributions
measured by SMT in living cells are not exponential but power-law. The
methodological insight is that H2B survival-time distributions in live
cells contain three components, of which the slowest reflects pure
photobleaching of the fluorophore rather than any chromatin-residence
process; dividing measured TF survival distributions by this H2B-derived
photobleaching kernel yields the photobleaching-corrected residence
distribution, which is heavy-tailed and best described by a continuum of
dissociation-rate states. This is a different and more direct route to
the same conclusion subsequently reached by GRID-style algorithms
\cite{Popp2021}. We adopt the Garcia \textit{et al.} result as the
empirical anchor for distinguishing
\emph{per-site} dwell times (exponential, with site-specific
$k_{\mathrm{off}}$) from the \emph{TF-perspective} residence-time
distribution (population-level mixture of per-site exponentials,
empirically heavy-tailed). The proofreading derivation in
Section~\ref{sec:gated} operates at the per-event level and so retains
its single-exponential per-engagement form. The implicit assumption ---
made explicit in Section~\ref{ssec:mixture} --- is that the long-residence
mass of the SMT distribution corresponds to TFs at cognate response
elements, and the short-residence mass to noncognate or transient
binding.

\paragraph{Mazzocca \textit{et al.} 2024~\cite{Mazzocca2024}.}
MINFLUX+HILO measurements of H2B dynamics under acute cohesin (RAD21)
depletion in mammalian interphase nuclei find a small effect on
single-locus diffusivity and no detectable shift in the anomalous
exponent~$\alpha$. This is the empirical signature predicted by
Shi~\textit{et al.}~\cite{Shi2025}: cohesin loss perturbs the global
$\mathbf{K}$ structure (centralities, eigenmode gap) while the local
fractional-Langevin exponent, set by short-time monomer dynamics, is
preserved. We use this dataset, together with pEM analysis of H2B in
3617 cells under RAD21 depletion (which shows shifted state fractions
without a detectable within-state~$\alpha$ change), as the calibration
anchor for the chromatin axis of~\eqref{eq:star} on the mouse side.

\paragraph{Lucas \textit{et al.} 2018~\cite{Lucas2018}.} MS2 imaging at
\textit{hb} establishes that the observed boundary precision exceeds simple
diffusion-limited Berg--Purcell predictions, motivating ``additional
mechanisms beyond simple diffusive readout.'' The Bcd-hub
work~\cite{Mir2017,Mir2018} provides one such mechanism (locally elevated
$c_{\mathrm{free}}$); the polymer-FPT framework here provides another
(short $\tau_c$ in pre-TAD chromatin), and the two are not mutually
exclusive.

\paragraph{Chen \textit{et al.} 2018~\cite{Chen2018}.} The first
simultaneous live-cell measurement of EP topology and transcriptional
activity in fly nc14 embryos, at the \textit{eve} locus with a reporter
integrated 142\,kb upstream. Three topological states are resolved:
$O_{\mathrm{off}}$ (open, distal, $\sim$87\% of nuclei, mean EP separation
730\,nm), $P_{\mathrm{off}}$ (paired but transcriptionally inactive,
$\sim$6\%, 385\,nm), and $P_{\mathrm{on}}$ (paired and active, $\sim$7\%,
331\,nm). Four findings revise the framework. First, the open-to-paired
transition rate ($f_1 = 0.017\text{/min}$, $\sim$1\,hr mean) is
\emph{slower} than the naive polymer first-passage prediction from MSD
scaling by approximately a factor of eight, indicating that an entropic
or architectural step beyond simple polymer encounter is required, with
the \textit{homie} insulator providing this step at \textit{eve}. Second,
sustained physical proximity is necessary for transcription: EP distance
declines continuously to $\sim$340\,nm and is held there before firing,
ruling out hit-and-run activation. Third, transcription itself stabilizes
the paired state: the active-paired dwell ($1/b_3\approx 90$\,min) is
roughly an order of magnitude longer than the inactive-paired dwell
($1/b_2\approx 7$\,min), implying a strong positive feedback between
firing and contact stability. Fourth, integration of a reporter at this
distance from \textit{eve} reduces endogenous \textit{eve} expression by
5--20\% per stripe and produces $\sim$5-fold increased adult abdominal
patterning defects, demonstrating that EP topology is a limiting resource
with phenotypic consequences. Chen \textit{et al.} therefore provides the
empirical foundation that simultaneously calibrates the polymer-FPT term
($\tau_c$ longer than the naive scaling), establishes the necessity of
architectural facilitation, and reveals a non-renewal feedback mechanism
in which firing itself extends~$\tau_d$.

\paragraph{Saiz, Martinovic \textit{et al.} 2026~\cite{Saiz2026}.} A direct
test of the structural-loop component of the polymer-FPT picture in the
mouse system. Auxin-degron depletion of CTCF in 2D mouse gastruloids
during the fate-decision window (48--72\,h, coincident with Wnt-induced
symmetry breaking) impairs gastruloid elongation but \emph{leaves cell
differentiation, ATAC-seq lineage composition, and RNA-seq trajectories
intact}. Reconstitution with an N-terminally truncated CTCF (incapable of
cohesin-mediated loop extrusion stalling) rescues early phenotypes;
gastruloids only collapse later (post-168\,h), revealing a separable role
for CTCF looping in post-gastrulation fine-tuning. Together with the
observation that acute degradation of cohesin (Rad21) preserves a
substantial fraction of chromatin
contacts~\cite{Rao2017,Schwarzer2017,Hsieh2022}, this argues against a
strict ``structural-loop $\to$ low $\tau_c$ $\to$ accurate fate readout''
chain in mouse gastrulation. The polymer-FPT term during the fate-decision
window is likely more permissive than the rigid-TAD interpretation would
suggest; CTCF-mediated looping enters as a fine-tuning constraint
\emph{after} the fate decision rather than as a rate-limiting input to it.

\paragraph{Xiao, Hafner, Boettiger 2021~\cite{Xiao2021}.} An equilibrium
model that resolves the apparent paradox of TAD-boundary disruptions
changing EP contact frequency by less than $\sim$2-fold but changing gene
expression by $\sim$10-fold, via a futile-cycle bistability with
effective Hill coefficient $n_H\!\gtrsim\!5$. We retain it as a
phenomenological description of how small structural perturbations
produce large transcriptional changes, but emphasize that the underlying
mechanism in living tissue is more naturally read off Chen \textit{et al.}
\cite{Chen2018}: the ten-fold longer dwell of the transcriptionally
active paired state, $P_{\mathrm{on}}$, relative to the inactive paired
state, $P_{\mathrm{off}}$, is a dynamic positive-feedback gain stage that
the equilibrium futile-cycle model approximates only in projection. The
practical consequence for the bound below is that firing itself extends
$\tau_d$, breaking the renewal-process assumption used in
\eqref{eq:Nprod} once the system has entered the active branch.

\paragraph{Kaizu \textit{et al.} 2014~\cite{Kaizu2014}.} ``The Berg--Purcell
limit revisited.'' Validates the original Berg--Purcell scaling against
single-receptor, finite-binding-kinetics simulations. Orthogonal to the
present work; cited to disambiguate from the title.

\section{The compound bound}
\label{sec:bound}

Consider a target gene whose transcriptional output during a window~$T$
requires three steps: (i)~TF occupancy of a cognate enhancer site, with
steady-state probability~$p_b(c)$ at TF free concentration~$c$; (ii)~EP
contact, occupying a fraction~$f_{\EP}$ of time with mean inter-contact
interval~$\tau_c$ and per-contact dwell~$\tau_d$; (iii)~successful firing
of a transcriptional burst, which requires the contact to persist for at
least a characteristic time~$\tau_{\mathrm{fire}}$ (the timescale of PIC
assembly and Pol II initiation, $\sim$tens of seconds). For an
approximately exponential dwell-time distribution with mean~$\tau_d$, the
per-contact firing probability is $p_{\mathrm{fire}} = 1 -
\exp(-\tau_d/\tau_{\mathrm{fire}})$ in the limit where $\tau_d$ is the
relevant cutoff. The number of productive contacts in window~$T$ is
\begin{equation}
  N_{\mathrm{prod}} \;=\; \frac{T}{\tau_c}\;p_{\mathrm{fire}}
  \;\approx\;
  \frac{T}{\tau_c}\;\frac{\tau_d}{\tau_{\mathrm{fire}}}
  \quad\text{(for }\tau_d \ll \tau_{\mathrm{fire}}\text{).}
  \label{eq:Nprod}
\end{equation}
The fractional variance of the gene's output decomposes as
\begin{equation}
  \left(\frac{\delta n}{\langle n\rangle}\right)^{\!2}
  \approx
  \underbrace{\left(\frac{\delta p_b}{p_b}\right)^{\!2}}_{\substack{\text{TF counting +}\\\text{specificity}}}
  +
  \underbrace{\frac{1}{N_{\mathrm{prod}}}}_{\substack{\text{polymer}\\\text{(both signs)}}}.
  \label{eq:decomp}
\end{equation}
The first term inherits both the Berg--Purcell diffusive-counting noise
\eqref{eq:bp} and the Lammers~\cite{Lammers2023} specificity penalty. The
second is the \emph{two-sided} polymer constraint: $1/N_{\mathrm{prod}}$
diverges either when $\tau_c$ is too large (the classical Berg--Purcell-on-EP
limit, ``not enough attempts'') or when $\tau_d$ is too small relative to
$\tau_{\mathrm{fire}}$ (the entropic limit, ``each attempt is too brief to
fire''). Because both~$\tau_c$ and~$\tau_d$ are set by the same polymer
dynamics, this term is not minimized by simply making the polymer faster;
fast polymers shrink~$\tau_d$ as well as~$\tau_c$, and the relevant
dimensionless quantity is the ratio $\tau_c/\tau_d$ multiplied by
$\tau_{\mathrm{fire}}/T$:
\begin{equation}
  \frac{1}{N_{\mathrm{prod}}}
  \;\approx\;
  \frac{\tau_c}{\tau_d}\cdot\frac{\tau_{\mathrm{fire}}}{T}.
  \label{eq:two_sided}
\end{equation}

The compound bound, applicable per locus per nucleus, is
\begin{empheq}[box=\fbox]{equation}
  \left(\frac{\delta c}{c}\right)^{\!2}_{\!\text{total}}
  \;\gtrsim\;
  \frac{1}{n_H^{2}}\!\left[
  \underbrace{\frac{1}{a\,D_{\eff}\,c_{\mathrm{free}}\,T}}_{\text{Berg--Purcell}}
  \;+\;
  \underbrace{\Phi\!\left(\frac{w/c}{\alpha_{Lammers}}\right)\frac{1}{s^{2}p^{2}\,T}}_{\text{Lammers specificity/sharpness}}
  \;+\;
  \underbrace{\frac{\tau_c(\alpha,s,b)}{\tau_d(\alpha,b)}\cdot\frac{\tau_{\mathrm{fire}}}{T}}_{\text{two-sided polymer}}
  \right].
  \label{eq:star}
\end{empheq}
The function~$\Phi$ encodes the scaling of the Lammers IR penalty with
noncognate interference; for low~$w/c$ it is order unity, and for high
$w/c$ it grows polynomially unless rescued by energy dissipation, in which
case sharpness, precision, and specificity face the tradeoff bound
analyzed in~\cite{Lammers2023}. The prefactor~$n_H$ is the effective Hill
coefficient of the gene-circuit input--output function. For a hypersensitive
readout in the Xiao futile-cycle regime~\cite{Xiao2021}, $n_H\gtrsim 5$ and
the apparent positional precision improves by~$1/n_H^{2}\gtrsim 1/25$,
provided the upstream noise terms reflect \emph{stochastic} fluctuations
about a fixed mean. The same gain amplifies systematic spatial gradients
in~$c$ and~$\tau_c$ at the same rate; the readout therefore relies on
mechanisms (multi-enhancer averaging, paracrine spatial averaging) that
suppress stochastic~$\delta\tau_c$ without diminishing the systematic
component. Each of the three terms scales as~$1/T$, but the prefactors are
set by physically independent processes: diffusive TF counting
(Berg--Purcell), gene-circuit kinetic architecture plus noncognate
interference (Lammers), and chromatin polymer dynamics (polymer FPT).

\section{Polymer first-passage scaling}
\label{sec:scaling}

For a chromatin locus undergoing fractional Langevin motion with monomer
MSD $\langle\Delta r^{2}(t)\rangle = \Gamma_{\alpha}\,t^{\alpha}$, the
mean first-passage time for two monomers at equilibrium 3D
separation~$R(s)$ to come within capture radius~$b$ obeys, in the
Amitai--Holcman / Gu\'erin--B\'enichou--Voituriez
form~\cite{Amitai2013,Guerin2016},
\begin{equation}
  \tau_c
  \;\sim\;
  R(s)^{2}\,\left(\frac{R(s)}{b}\right)^{2/\alpha-1}\,
  \Gamma_{\alpha}^{-1/\alpha},\qquad \alpha\le 1.
  \label{eq:tau_c}
\end{equation}
Equation~\eqref{eq:tau_c} should be read as a \emph{homopolymer reference}
rather than as the operative law for chromatin. Shi, Shin, and
Thirumalai~\cite{Shi2025} show that an effective-homopolymer reduction of
the wild-type chromatin distance map gives the wrong exponents
($\tau\sim s^{1.1}$ and $\tau\sim\langle r\rangle^{4}$, vs.\ the
experimentally measured $\tau\sim s^{0.7\text{--}0.8}$ and
$\tau\sim\langle r\rangle^{2.6}$), and randomly shuffling the WT contact
map at fixed polymer connectivity destroys the anomalous scaling
entirely. The operative quantity is therefore not~\eqref{eq:tau_c} but the
relaxation time and first-passage time computed from the connectivity
matrix~$\mathbf{K}$ derived from the static contact map (HIPPS-DIMES). For
single first-passage to a contact threshold~$r_c$ they obtain the
Szabo--Schulten--Schulten-like result $\tau_c\propto
D^{-1}r_c^{-1}\langle r\rangle^{3}$, in quantitative agreement with both
their theoretical and Br\"uckner's experimental
trajectories~\cite{Bruckner2023,Shi2025}. We retain~\eqref{eq:tau_c}
below for dimensional and intuitional purposes only; the operative
$\tau_c$ used in~\eqref{eq:star} is the empirically measured (or
HIPPS-DIMES-computed) value at the readout locus.

\begin{table}[h]
\centering
\begin{tabular}{lccc}
\toprule
Regime & $\alpha$ & $2/\alpha-1$ & $\tau_c$ at $s\!\sim\!100$\,kb \\
\midrule
Rouse / phantom chain          & 0.5  & 3   & seconds--tens of seconds \\
Crumpled / mild crowding       & 0.4  & 4   & minutes \\
Heterochromatic / TAD-confined & 0.25--0.30 & 5.7--7 & tens of minutes--hours \\
\bottomrule
\end{tabular}
\caption{Homopolymer reference values for first-passage scaling. The
exponent $2/\alpha-1$ in~\eqref{eq:tau_c} drives $\tau_c$ over several
decades as~$\alpha$ varies in the biologically relevant range.
\textbf{These regimes do not classify chromatin loci}: Shi
\textit{et al.}~\cite{Shi2025} demonstrate that the WT chromatin contact
map produces dynamics not captured by any single~$\alpha$ (the
connectivity-matrix eigenvalue spectrum exhibits three distinct scaling
regimes, $|\lambda_p|\propto p^{1.2},\,p^{3},\,p^{1.5}$, across mode
indices). The table is included for orientation against the homopolymer
literature; the operative~$\tau_c$ is locus-specific and set by
$\mathbf{K}$.}
\label{tab:regimes}
\end{table}

\paragraph{Empirical correction and structural reframing.} Br\"uckner
\textit{et al.}~\cite{Bruckner2023} report that EP encounter times in fly
nuclei depend on genomic separation \emph{less strongly}
than~\eqref{eq:tau_c} would predict. Shi, Shin, and
Thirumalai~\cite{Shi2025} resolve this as a structural rather than
empirical correction: the rapid dynamics is encoded in the WT contact
map's specific pattern of long-range constraints (likely mediated by
cohesin and transcription factors) and is reproduced quantitatively by
the HIPPS-DIMES connectivity matrix using only static Micro-C input. The
practical effect is to push the operative~$\tau_c$ \emph{below} the
homopolymer reference of Table~\ref{tab:regimes}, favoring compatibility
with developmental decision windows; but the framework should be read
with $\tau_c$ as the contact-map-derived (or directly measured) quantity,
not as a function of $(\alpha,s,b)$.

\paragraph{The dwell time and the two-sided constraint.} In the homopolymer
reference, the same parameters that set $\tau_c$ also set the per-contact
dwell time $\tau_d \sim b^{2}\,\Gamma_{\alpha}^{-1/\alpha}$, and the ratio
$\tau_c/\tau_d \sim (R/b)^{2/\alpha-1}$ depends only on the geometric ratio
$R/b$ and the exponent. For fly EP separations ($R\!\sim\!200$\,nm,
$b\!\sim\!50$\,nm), this gives $\tau_c/\tau_d \sim (R/b)^{3} \approx 64$
in the Rouse regime: many attempts, each short. With
$\tau_{\mathrm{fire}} \sim 10\text{--}30$\,s and bare-polymer
$\tau_d \lesssim 1$\,s, individual contacts fire with low probability,
consistent with the bursty, low-duty-cycle transcription observed at
\textit{eve} stripes~\cite{Bothma2014,Lammers2020}. Under
HIPPS-DIMES~\cite{Shi2025} both $\tau_c$ and the relaxation time $\tau$
are set by eigenmodes of $\mathbf{K}$, so the closed-form ratio above is
not the operative quantity for chromatin; we retain the estimate as a
benchmark against which calibrated values from Chen
\textit{et al.}~\cite{Chen2018} and Mazzocca
\textit{et al.}~\cite{Mazzocca2024} are compared.

The Br\"uckner result~\cite{Bruckner2023} is therefore especially
informative: their measured EP contacts are not just more frequent than
classical scaling would predict but also \emph{persistently associated}
along the genomic distance axis, suggesting non-Markovian polymer
dynamics that selectively extend $\tau_d$ relative to the simple
fractional-Langevin prediction. This is the empirical resolution of the
fly side of~\eqref{eq:two_sided}: short~$\tau_c$ \emph{and}, anomalously,
non-vanishing~$\tau_d$.

\paragraph{Empirical calibration in fly nc14 (Chen 2018).} Chen
\textit{et al.}~\cite{Chen2018} provide direct measurements of all three
relevant timescales at \textit{eve} with an EP separation of 142\,kb. The
mean open-to-paired transition time is $\tau_c \approx 1$\,hr, eight times
longer than the naive polymer first-passage prediction; the inactive-paired
dwell is $\tau_d^{\mathrm{off}}\approx 7$\,min; the active-paired dwell is
$\tau_d^{\mathrm{on}}\approx 90$\,min. Two consequences follow. (i)~The
naive scaling~\eqref{eq:tau_c} \emph{underestimates}~$\tau_c$ in fly nuclei,
opposite in sign to the Br\"uckner long-range-correlation
correction~\cite{Bruckner2023}; the two effects coexist at different EP
separations and the empirically relevant~$\tau_c$ depends on locus identity
and architectural elements. (ii)~Once a contact has fired, the relevant
dwell jumps from $\tau_d^{\mathrm{off}}$ to $\tau_d^{\mathrm{on}}$ ---
roughly thirteen-fold --- so the bound~\eqref{eq:two_sided} should be read
with $\tau_d = \tau_d^{\mathrm{on}}$ for productive bursts. The
positive-feedback mechanism implies that the renewal approximation
underlying \eqref{eq:Nprod} fails in the active branch, where successive
contacts are correlated through the persistence of the $P_{\mathrm{on}}$
state.

\paragraph{Architectural and transcription-stabilized dwell extension.}
Both fly and mouse readouts circumvent the bare-polymer dwell-time
limitation through dynamic mechanisms that extend~$\tau_d$ once contact
is achieved, but via parallel and physically distinct routes. In fly
nc14, insulator pairing (e.g.\ \textit{homie--homie}~\cite{Chen2018})
brings the loci into contact, and \emph{transcription itself} stabilizes
the paired state by an order of magnitude (the
$P_{\mathrm{off}}\!\to\!P_{\mathrm{on}}$ transition with
$\tau_d^{\mathrm{on}}\!\sim\!90$\,min vs.\ $\tau_d^{\mathrm{off}}\!\sim\!7$\,min,
as measured by Chen \textit{et al.}~\cite{Chen2018}). In mouse, cohesin-
and CTCF-anchored loops~\cite{Gabriele2022,Mach2022} are metastable
trapped states with lifetimes of tens of minutes,
$\tau_d^{\mathrm{loop}} \gtrsim \SI{e3}{\second} \gg \tau_{\mathrm{fire}}$,
occupied only a few percent of the cell cycle. In both cases the system
trades a high-rate / low-success regime for a low-rate / high-success
regime, improving $N_{\mathrm{prod}}$ when $\tau_d$ exceeds the firing
threshold. The fly mechanism is locus-encoded (insulator sequences plus
transcription-coupled stabilization); the mouse mechanism is genome-wide
(CTCF/cohesin extrusion). They are not rival hypotheses but parallel
implementations of the same physical strategy.

This reframing is consistent with the Saiz \textit{et al.} CTCF-AID
gastruloid result~\cite{Saiz2026}: during fate specification, the system
operates by sorting between bistable states (Sox2-high vs.\ Bra-high),
which tolerates short~$\tau_d$ because the multistable dynamics lock in
the decision once made. After fate specification, accurate expression
levels require many productive bursts per unit time, which requires
$\tau_d \gtrsim \tau_{\mathrm{fire}}$, which requires CTCF-anchored loops
--- recovering the observed late-development requirement for CTCF's
N-terminal extrusion-stalling function.

\paragraph{A polymer-derived Hill law.} Shi, Shin, and
Thirumalai~\cite{Shi2025} derive an explicit Hill-form transcription rate
from the polymer first-passage statistics: assuming the rate is
$k = 1/(\tau_c+\tau_d)$ and that $\tau_c$ scales with contact probability
as $\tau_c\sim\tau_0\,p_c^{-\theta}$, the normalized rate becomes
\begin{equation}
  \hat k \;=\; \frac{k}{k_{\max}}
  \;=\; \frac{1}{1 + (\tau_0/\tau_d)\,p_c^{-\theta}},
  \label{eq:polymer_hill}
\end{equation}
with effective ``cooperativity''~$\theta$. This supplies a mechanistic
bridge from~$\tau_c$ to the prefactor~$n_H$ in~\eqref{eq:star} that is
\emph{independent of and complementary to} the
Xiao--Hafner--Boettiger~\cite{Xiao2021} equilibrium futile-cycle
amplifier: small structural perturbations to $p_c$ produce large changes
in $\hat k$ through the polymer-FPT response itself, without requiring
the equilibrium gene-circuit nonlinearity. Where the latter is invoked in
the Discussion, the polymer-FPT Hill law operates upstream and stacks
with it.

\paragraph{Compatibility window for fly nc14.} Saturating the Gregor 10\%
positional precision target with $T = \SI{180}{\second}$ requires
$\tau_c \lesssim \SI{2}{\second}$ at the single-locus level, loosening to
$\tau_c \lesssim \SI{10}{\second}$ with Gregor's spatial averaging over
$\sim$5 syncytial neighbors. Only the Rouse-like regime
($\alpha\!\approx\!0.5$) reaches this window without the
Br\"uckner correction; with the correction, the
crumpled/$\alpha\!\approx\!0.4$ regime becomes marginally accessible.

\section{Polymer-gated kinetic proofreading: a formal coupling}
\label{sec:gated}

The compound bound \eqref{eq:star} writes the polymer and Lammers terms
additively, which captures their dimensional independence but not the
dynamical coupling between them. Kim \& Larson~\cite{KimLarson2025} have
recently extended the Gregor timing framework with an explicit
kinetic-proofreading step, arriving at a speed--efficiency--specificity
tradeoff $\psi\approx 1-E/M$ for transcriptional bursting in development.
Their treatment, like Lammers~\cite{Lammers2023}, takes the substrate
(transcription factor) as available at the locus and the proofreading
cycle as running freely once initiated; the polymer step that brings
enhancer and promoter into productive contact does not enter. We close
that gap here. The Lammers framework treats decision-time accrual and
kinetic-proofreading specificity gain as properties of a gene-circuit
transition matrix~$\mathbf{K}_{\text{Lammers}}$ whose dynamics presumes
the gene to be in burst-cycling mode throughout the decision window. Each proofreading cycle requires the locus to
traverse $N_A$ activation steps in the active conformation, which in turn
requires sustained EP contact. Finite contact dwell times $\tau_d$
therefore throttle the rate at which proofreading cycles complete and
degrade the realized specificity. We derive the coupling here and show
that, in contrast to the naive duty-cycle suppression
$\mathrm{IR}_{\mathrm{eff}} = f_{\EP}\cdot\mathrm{IR}_{\mathrm{Lammers}}$,
the coupled treatment produces a \emph{threshold} at $\tau_d\sim
N_A\,\tau_{\mathrm{step}}$ rather than a smooth proportional reduction.

\subsection{Gated cycle dynamics}

Let $\tau_{\mathrm{step}}$ denote the mean time per Lammers activation
step (binding-site engagement, Mediator engagement, PIC assembly,
Pol\,II loading; $\tau_{\mathrm{step}}\sim 10$--$100$\,s for typical
chromatin steps). Within a single EP engagement of mean duration
$\tau_d$, model both forward step completion (rate
$1/\tau_{\mathrm{step}}$) and disengagement (rate $1/\tau_d$) as competing
Poisson processes; reverse steps are neglected, consistent with the
Hopfield--Ninio assumption of irreversibility for the
specificity-amplifying transitions. The probability of completing one
activation step before disengagement is
\begin{equation}
  p_{\mathrm{step}}
  \;=\; \frac{1/\tau_{\mathrm{step}}}{1/\tau_{\mathrm{step}}+1/\tau_d}
  \;=\; \frac{\tau_d}{\tau_d + \tau_{\mathrm{step}}}.
  \label{eq:p_step}
\end{equation}
The probability that a given engagement completes the full $N_A$-step
cycle is therefore
\begin{equation}
  P_{\mathrm{cycle}}(\tau_d;\,N_A,\tau_{\mathrm{step}})
  \;=\; p_{\mathrm{step}}^{\,N_A}
  \;=\; \left(\frac{\tau_d}{\tau_d + \tau_{\mathrm{step}}}\right)^{\!N_A}.
  \label{eq:P_cycle}
\end{equation}
This expression has two limits. For $\tau_d \gg N_A\,\tau_{\mathrm{step}}$,
$P_{\mathrm{cycle}}\to 1$: every engagement realizes a full proofreading
cycle and the Lammers framework applies as written. For $\tau_d \ll
\tau_{\mathrm{step}}$, $P_{\mathrm{cycle}}\sim(\tau_d/\tau_{\mathrm{step}})^{N_A}$:
no engagement realizes a cycle and the gene operates without proofreading
specificity. The transition between these limits is sharp; for $\tau_d =
N_A\,\tau_{\mathrm{step}}$, $P_{\mathrm{cycle}}\to 1/e$ in the large-$N_A$
limit, locating the half-completion threshold at $\tau_d^{\mathrm{thresh}}
\approx N_A\,\tau_{\mathrm{step}}$.

\subsection{Realized specificity}

Let each completed activation step impose a multiplicative discrimination
factor~$\alpha < 1$ on noncognate vs.\ cognate firings (the
Hopfield--Ninio amplifier). If engagements that fall off after $n<N_A$
steps reset the proofreading chain (cognate and noncognate both equally
discarded), the realized specificity averaged over engagements is
\begin{equation}
  f_{\mathrm{realized}}^{-1}
  \;=\; \mathbb{E}\!\left[\alpha^{-(1+N_c)}\right]
  \;=\; (1-p_{\mathrm{step}})\sum_{n=0}^{N_A-1} p_{\mathrm{step}}^{\,n}\,\alpha^{-(1+n)}
  \;+\; p_{\mathrm{step}}^{\,N_A}\,\alpha^{-(N_A+1)},
  \label{eq:f_realized}
\end{equation}
which interpolates between $f_{\mathrm{realized}}=\alpha$ for
$\tau_d\to 0$ (no proofreading) and $f_{\mathrm{realized}}=\alpha^{N_A+1}$
for $\tau_d\to\infty$ (full Lammers ceiling).

\subsection{Coupled information rate}

Engagements occur at rate $(\tau_c+\tau_d)^{-1}$. The rate of
\emph{completed} proofreading cycles is
\begin{equation}
  \mathcal{R}_{\mathrm{cycle}}
  \;=\; \frac{P_{\mathrm{cycle}}(\tau_d;N_A,\tau_{\mathrm{step}})}{\tau_c + \tau_d}.
  \label{eq:R_cycle}
\end{equation}
In the high-noncognate-interference regime ($w/c \gg \alpha_{\mathrm{L}}$,
mouse-relevant), the per-cycle information accrual is dominated by the
specificity-controlled noise floor and scales as $f_{\mathrm{realized}}^{2}$
\cite{Lammers2023}. The coupled effective information rate is therefore
\begin{empheq}[box=\fbox]{equation}
  \mathrm{IR}_{\mathrm{eff}}(\tau_c,\tau_d;N_A,\tau_{\mathrm{step}})
  \;\propto\;
  \underbrace{\frac{1}{\tau_c+\tau_d}}_{\text{engagement rate}}
  \;\cdot\;
  \underbrace{\left(\frac{\tau_d}{\tau_d+\tau_{\mathrm{step}}}\right)^{\!N_A}}_{P_{\mathrm{cycle}}}
  \;\cdot\;
  \underbrace{f_{\mathrm{realized}}^{\,2}(\tau_d;N_A,\alpha)}_{\text{IR per completed cycle}}.
  \label{eq:IR_eff}
\end{empheq}
The naive worst-case treatment, $\mathrm{IR}_{\mathrm{eff}}=f_{\EP}\,
\mathrm{IR}_{\mathrm{Lammers}}$ with $f_{\EP}=\tau_d/(\tau_c+\tau_d)$,
differs from \eqref{eq:IR_eff} in the second and third factors. Both
naive and coupled expressions agree in the saturated regime
$\tau_d\gg N_A\tau_{\mathrm{step}}$, where $P_{\mathrm{cycle}}\to 1$ and
$f_{\mathrm{realized}}\to\alpha^{N_A+1}$. Below saturation, the coupled
expression collapses as $\tau_d^{N_A}$ rather than as $\tau_d$, a
qualitatively different scaling.

\subsection{Optimal dwell time}

At fixed $\tau_c$ and in the saturated-specificity regime
($f_{\mathrm{realized}}\approx\alpha^{N_A+1}$, independent of $\tau_d$),
maximizing $\mathrm{IR}_{\mathrm{eff}}$ over $\tau_d$ via $\partial_{\tau_d}
\ln\mathrm{IR}_{\mathrm{eff}}=0$ yields
\begin{equation}
  \tau_d^{*}
  \;=\; \frac{\tau_{\mathrm{step}}(N_A-1) +
        \sqrt{\tau_{\mathrm{step}}^{2}(N_A-1)^{2} + 4N_A\,\tau_{\mathrm{step}}\,\tau_c}}{2}.
  \label{eq:tau_d_opt}
\end{equation}
Two limits clarify the structure. For rare engagements
($\tau_c \gg N_A\tau_{\mathrm{step}}$, mouse), $\tau_d^{*}\approx
\sqrt{N_A\,\tau_{\mathrm{step}}\,\tau_c}$: the optimal dwell is the
geometric mean of the per-step time and the inter-contact interval. For
frequent engagements ($\tau_c \ll N_A\tau_{\mathrm{step}}$, fly nc14
small-$s$ regime under the Br\"uckner correction~\cite{Bruckner2023}),
$\tau_d^{*}\approx (N_A-1)\,\tau_{\mathrm{step}}$: the optimal dwell is
just long enough to complete one cycle. In both regimes,
$\tau_d^{*}\gtrsim N_A\,\tau_{\mathrm{step}}$, confirming the threshold
identified in~\eqref{eq:P_cycle}.

\subsection{Mixture distributions and the cognate-site assumption}
\label{ssec:mixture}

The derivation above treats the contact dwell $\tau_d$ as exponentially
distributed with a single mean. SMT measurements of transcription-factor
residence times are not exponential at the population level: Garcia,
Upadhyaya, Hager and colleagues~\cite{Garcia2021NAR}, by dividing
measured TF survival distributions by an H2B-derived photobleaching
kernel, showed that the photobleaching-corrected population
residence-time distribution in living cells is power-law and is best
described by a continuum of dissociation-rate states rather than 2--3
discrete components. Two distributions
must be distinguished. The \emph{per-site} dwell time at any single
binding site is well approximated as a single exponential set by the
site's local off-rate. The \emph{TF-perspective} residence-time
distribution measured by SMT is a mixture of these per-site exponentials
weighted by visitation, and is heavy-tailed at the population level even
when each per-site dwell is exponential.

The Hopfield--Ninio proofreading cycle runs to completion within one
engagement at one site, so the per-event ceiling
$f\to\alpha^{N_A+1}$ is set by the per-site exponential, and the
derivation of \eqref{eq:P_cycle}--\eqref{eq:IR_eff} is unchanged at the
single-event level. The implicit assumption underlying
Section~\ref{sec:gated} is therefore that \emph{the long-residence
component of the SMT distribution corresponds to TFs bound at cognate
response elements}, while the short-residence mass reflects noncognate
or transient binding. This identification is consistent with the
biological expectation that high-affinity cognate sites have low
$k_{\mathrm{off}}$ and with empirical findings that long-lived dwell-time
states are enriched at functionally engaged target loci~\cite{Garcia2021NAR}.

The EP-contact dwell $\tau_d$ that enters \eqref{eq:IR_eff} is itself a
two-component mixture in the empirical data: a short bare-polymer
component and a long architecturally-stabilized component
(cohesin/CTCF-anchored loops in mouse~\cite{Gabriele2022,Mach2022};
insulator-paired states with transcription-coupled stabilization in fly
\cite{Chen2018}). With long-fraction $\phi$, long-component mean
$\tau_d^{L}$, and short-component mean $\tau_d^{S}$, the cycle-completion
probability \eqref{eq:P_cycle} averages to
\begin{equation}
  P_{\mathrm{cycle}}^{\mathrm{mix}}
  \;=\; (1-\phi)\!\left(\frac{\tau_d^{S}}{\tau_d^{S}+\tau_{\mathrm{step}}}\right)^{\!N_A}
  + \phi\!\left(\frac{\tau_d^{L}}{\tau_d^{L}+\tau_{\mathrm{step}}}\right)^{\!N_A}.
  \label{eq:P_cycle_mix}
\end{equation}
For mouse parameters ($\phi\sim 0.05$, $\tau_d^{L}\sim 10^{3}$\,s,
$\tau_d^{S}\sim 1$\,s, $N_A = 4$, $\tau_{\mathrm{step}}\sim 10$\,s), the
short component contributes $\sim 7\times 10^{-5}$ and the long
component contributes $\sim 0.96$, giving
$P_{\mathrm{cycle}}^{\mathrm{mix}}\approx\phi$. The threshold result of
Section~\ref{sec:gated} is therefore better stated as: \emph{a small
fraction $\phi$ of EP contacts (the long-component) carries essentially
all the proofreading}, with the operative quantity being
$\phi\,\tau_d^{L}$ relative to $N_A\tau_{\mathrm{step}}$ rather than the
mean $\langle\tau_d\rangle$. The biological mechanisms that we have
described as ``extending $\tau_d$'' --- CTCF/cohesin loops, insulator
pairing, transcription-coupled stabilization --- are mechanisms for
\emph{generating the long-component fraction $\phi$} on top of an
otherwise short bare-polymer distribution; they do not shift the bulk
of the distribution, only the tail.

The Berg--Purcell term in \eqref{eq:star} is, by the same distinction,
governed by the per-site exponential at the cognate enhancer and
retains its canonical $\delta c/c \sim T^{-1/2}$ scaling. Heavy-tailed
TF-perspective residence enters BP only through two upstream channels:
reduction of effective $c_{\mathrm{free}}$ by noncognate sequestration,
and bursty, non-Poisson arrival statistics at the cognate site that
violate the renewal assumption underlying~\eqref{eq:Nprod}. Both effects
bias the Lammers $w/c$ ratio toward heavy-tail-fluctuating values rather
than a single number, and should be read as additional --- not
replacement --- corrections to the bound.

\subsection{Implications}

Three consequences depart materially from the additive treatment.

\paragraph{Threshold rather than smooth suppression.} The naive
$f_{\EP}$-suppression predicts that any reduction in $\tau_d$ degrades the
Lammers term proportionally. The coupled treatment \eqref{eq:IR_eff}
predicts a sharp threshold at $\tau_d \approx N_A\,\tau_{\mathrm{step}}$:
above it, full Lammers performance is realized; below it, performance
collapses as $(\tau_d/\tau_{\mathrm{step}})^{N_A}$. For $N_A\!\sim\!4$
(mouse-relevant per~\cite{Lammers2023}), this is a $\tau_d^{4}$ collapse
--- four orders of magnitude steeper than the linear duty-cycle scaling.
Operating windows that look marginal under the additive bound are either
comfortably inside or catastrophically outside the coupled bound.

\paragraph{Mouse requires a long-component fraction $\phi$ specifically.}
With $\tau_{\mathrm{step}}\sim 10$--$100$\,s and $N_A\sim 4$, the
single-component threshold $N_A\tau_{\mathrm{step}}$ falls in the range
$40$--$400$\,s. The empirical EP-contact dwell distribution in mouse is
bimodal, with $\tau_d^{S}\sim 1$\,s (bare polymer) and $\tau_d^{L}\gtrsim
10^{3}$\,s (cohesin/CTCF-anchored loops) at long-component fraction
$\phi\sim$ a few percent~\cite{Gabriele2022,Mach2022}. By
\eqref{eq:P_cycle_mix}, cycle completion is essentially carried by the
long component: $P_{\mathrm{cycle}}^{\mathrm{mix}}\approx\phi$ when
$\tau_d^{L}\gg N_A\tau_{\mathrm{step}}\gg\tau_d^{S}$. The coupled
framework therefore identifies a quantitative reason --- beyond the
appeal to enhancer redundancy --- why mouse \emph{requires} cohesin-
anchored loops to make Lammers proofreading viable: not because the
loops shift the average contact dwell, but because they generate the
long-component tail $\phi$ that does almost all the proofreading work.
Eliminating the loop population (acute cohesin/CTCF
depletion~\cite{Saiz2026}) collapses $\phi\to 0$ and drives
$P_{\mathrm{cycle}}^{\mathrm{mix}}\to (\tau_d^{S}/(\tau_d^{S}+
\tau_{\mathrm{step}}))^{N_A}\sim 10^{-5}$, even though the mean contact
frequency~\cite{Hsieh2022} is largely preserved. Conversely, fly nc14
with $N_A\!\sim\!1$ has single-component threshold $\sim 10^{1}$\,s,
already exceeded by the inactive-paired dwell
$\tau_d^{\mathrm{off}}\sim 7$\,min reported by Chen
\textit{et al.}~\cite{Chen2018} for the dominant component of the
distribution; fly does not need a separate long-component to exceed
threshold.

\paragraph{Resolution of the Saiz fate-vs-morphogenesis dissociation.}
Fate decisions are bistable threshold-crossings tolerant of
$f_{\mathrm{realized}}$ values well below the Lammers ceiling: any
$f_{\mathrm{realized}}>\alpha$ together with downstream multistability
locks in the decision. Fate-relevant proofreading therefore runs at
modest effective $N_A^{\mathrm{fate}}\sim 1$--$2$, with threshold
$\tau_d^{\mathrm{thresh}}\sim\tau_{\mathrm{step}}\sim 10^{1}$\,s ---
accessible to bare-polymer dwell times. Morphogenesis depends on
quantitatively accurate transcriptional output across cells and therefore
on full-$N_A$ proofreading, with threshold
$\tau_d^{\mathrm{thresh}}\sim 10^{2}$--$10^{3}$\,s, accessible only with
cohesin-anchored loops. The Saiz \textit{et al.}~\cite{Saiz2026}
CTCF-AID phenotype --- fate intact, morphogenesis impaired --- is then
the predicted signature of perturbing $\tau_d$ across the
morphogenesis-relevant threshold while leaving the fate-relevant
threshold untouched. The Shi \textit{et al.}~\cite{Shi2025} signed
$\tau_{ij}$-shift result --- intra-TAD pairs slow down, boundary pairs
speed up under cohesin loss --- gives the locus-by-locus structure of
this perturbation: morphogenesis-critical genes whose enhancers sit
inside TADs are preferentially knocked across the threshold; boundary
genes are not. This is a sharper prediction than ``loops are dispensable
for fate.''

\paragraph{Energy efficiency and aborted cycles.} Each completed
activation step dissipates $\sim k_BT\ln(\eta_{ab}\eta_{ua}/\eta_{ba}\eta_{ub})$
of free energy~\cite{Lammers2023}. Aborted cycles dissipate only the
energy of steps actually completed, not the full $N_A$-step cost. The
information-per-energy ratio in the coupled framework is therefore
\begin{equation}
  \frac{\mathrm{IR}_{\mathrm{eff}}}{\Phi_{\mathrm{eff}}}
  \;\sim\;
  \frac{f_{\mathrm{realized}}^{\,2}\,p_{\mathrm{step}}^{N_A}}
       {\phi_{\mathrm{step}}\,p_{\mathrm{step}}\,(1-p_{\mathrm{step}}^{N_A})/(1-p_{\mathrm{step}})},
  \label{eq:efficiency}
\end{equation}
which recovers the Lammers efficiency in the saturated regime
$\tau_d\gg N_A\tau_{\mathrm{step}}$, but for $\tau_d\sim
\tau_{\mathrm{step}}$ exhibits a regime where partial cycles dissipate
without producing information. This identifies a previously
uncharacterized failure mode of biological kinetic proofreading: a
locus with a marginal contact dwell time pays the metabolic cost of
attempted proofreading without realizing the specificity benefit. The
prediction is testable through paired measurements of nascent
transcription burstiness and ATP consumption across acute cohesin
depletion conditions.

\paragraph{Caveats.} The derivation assumes (i) exponentially distributed
contact dwell times, (ii) irreversible forward activation steps, (iii)
that disengagement resets the proofreading chain. Reversible steps and
partial chain memory would soften the threshold; non-exponential
$\tau_d$ distributions (e.g., the bimodal active vs.\ inactive paired
states of Chen \textit{et al.}~\cite{Chen2018}) require evaluating
\eqref{eq:IR_eff} as an average over the dwell-time distribution and
would generally favor the long-dwell tail. The coupling derived here
should therefore be read as the leading-order treatment of an
unavoidable interaction between two physically independent constraints,
not as a final closed form.

\section{Multi-axis fly-vs-mouse asymmetry}
\label{sec:fly_vs_mouse}

The fly-vs-mouse asymmetry, often discussed as a single tradeoff between
equilibrium and nonequilibrium gene-circuit operation, in fact runs along
four physically independent axes that the foregoing sections have
introduced piecemeal. We collect them here and observe that
Lammers~\cite{Lammers2023} and Kim \& Larson~\cite{KimLarson2025}
implicitly conflate the first two.

\paragraph{Axis~1: species discrimination at the cognate site.} Set by
the noncognate-to-cognate \emph{species} concentration ratio $w/c$
\cite{Lammers2023} (or equivalently $M$ in Larson's
notation~\cite{KimLarson2025}). This is the question of how reliably a
gene locus, when occupied, identifies which TF species is bound. Resolved
by kinetic proofreading. Fly nc14 has $w/c\sim 47$ (low; equilibrium
suffices); mouse early embryo has $w/c\sim 10^{3}$ (high; proofreading
required).

\paragraph{Axis~2: site discrimination by the TF.} Set by the TF-perspective
ratio of noncognate to cognate sites in the genome,
$\rho \equiv N_{\mathrm{noncog}}/N_{\mathrm{cog}}$. For Bicoid in fly,
$\rho\sim 10^{2}$--$10^{3}$ (counting weak motif matches); for typical
mammalian TFs, $\rho\sim 10^{4}$--$10^{6}$, both because the genome is
larger and because mammalian regulatory landscapes are more permissive.
Crucially, this axis is \emph{not} captured by $w/c$: a low-$w/c$ system
(few competing species) still has high~$\rho$ if any single species must
discriminate among many candidate sites. Larson~\cite{KimLarson2025}
and Lammers~\cite{Lammers2023} elide this distinction by treating
``few privileged TFs in fly nc14'' as a complete answer; it is in fact
only an answer to Axis~1. The Axis~2 problem is real in fly and is
resolved by mechanisms orthogonal to proofreading: facilitated diffusion,
1D sliding, transcription-factor hubs/clusters that locally elevate
$c_{\mathrm{free}}$ at active enhancers (Bcd hubs in
fly~\cite{Mir2017,Mir2018}; Sox2/Bra clusters in mouse), and power-law
residence-time distributions with cognate-enriched long
tails~\cite{Garcia2021NAR}. In Eq.~\eqref{eq:star} this axis enters
implicitly through the $c_{\mathrm{free}}$ and $D_{\mathrm{eff}}$ that
appear in the Berg--Purcell prefactor; both are reduced relative to bulk
values by noncognate sequestration~\cite{Mirny2009} and locally elevated
by hub/cluster mechanisms.

\paragraph{Axis~3: enhancer-to-promoter contact formation.} Set by
chromatin polymer first-passage; entered as the $\tau_c/\tau_d$ term in
Eq.~\eqref{eq:star} and elaborated in Sections~\ref{sec:scaling}--\ref{sec:gated}.
Fly nc14 has short $\tau_c$ (Br\"uckner-corrected pre-TAD chromatin);
mouse has long $\tau_c$, partly rescued by cohesin loops. Operative
parameter: not the chromatin MSD exponent~$\alpha$ as a single number,
but the connectivity matrix~$\mathbf{K}$ derived from the static contact
map~\cite{Shi2025}.

\paragraph{Axis~4: cycle-completion threshold.} Set by the long-component
fraction $\phi$ of the EP-contact dwell distribution and the per-step
activation time $\tau_{\mathrm{step}}$, via the coupled bound
\eqref{eq:P_cycle_mix}. This axis is the mechanism by which Axes~1 and~3
interact: even a system that has resolved both species discrimination
($w/c$ low or proofreading available) and site discrimination
($c_{\mathrm{free}}$ adequate locally) fails if the proofreading cycle
cannot complete within a contact dwell. Mouse satisfies it through the
loop-anchored long-component $\phi\sim 0.05$ generated by CTCF/cohesin;
fly satisfies it through transcription-stabilized pairing
\cite{Chen2018} and the small effective $N_A$ of equilibrium circuits.

The four axes are dimensionally independent. A system can in principle
sit comfortably on any subset and at the wire on the rest; empirically,
both fly nc14 and mouse early embryo sit near the bound on \emph{all
four} simultaneously, which is why morphogen-based developmental
decisions are restricted to specific time windows and break down
outside them.

Lammers \textit{et al.}~\cite{Lammers2023} identify Axis~1 as the
parameter that distinguishes fly (equilibrium-permissive) from mouse
(energy-dissipation-required). Equation~\eqref{eq:star} adds the polymer
axes; Section~\ref{sec:gated} couples Axes~1 and~4. The remainder of this
section presents the joint Axis~1 + Axis~3 picture in tabular form,
with Axes~2 and~4 entering as the prefactor on Berg--Purcell
(noncognate sequestration / hub amplification) and the mixture-fraction
$\phi$ on the polymer term respectively. We caution, following Shi
\textit{et al.}~\cite{Shi2025}, that ``$\alpha$ at the readout locus'' is
not a clean per-system number: the distribution of single-locus~$\alpha$
within a single nucleus is broad (Shi
\textit{et al.}~\cite{Shi2025}, Fig.~6, $\langle\alpha\rangle\approx 0.52$
with substantial spread), and the per-locus diffusivity is anticorrelated
with a graph-theoretic centrality $C_i = \sum_{j\neq i} r_{ij}^{-m}$
computed from the contact map. The operative chromatin-axis parameter is
therefore better stated as the local structure of the connectivity
matrix~$\mathbf{K}$ at the readout locus, with $C_i$ and the eigenmode
contributions as measurable proxies. With that caveat, each embryo lives
at a point in the $(w/c,\,\text{chromatin axis})$ plane, and the location
determines which constraints are tight and which mechanisms must be in
place to satisfy them.

\begin{table}[h]
\centering
\small
\begin{tabular}{p{0.18\linewidth}cc p{0.30\linewidth} p{0.28\linewidth}}
\toprule
System & $w/c$ & $\alpha$ at readout locus & Lammers axis (kinetic) & Polymer-FPT axis (structural) \\
\midrule
\textit{Drosophila} nc14 &
$\sim 47$ &
$\sim 0.5$ (pre-TAD~\cite{Hug2017,Ogiyama2018}) &
Equilibrium gene circuits suffice (3--5 binding sites) &
Single-locus $\tau_c$ short; Bcd hubs~\cite{Mir2017,Mir2018} amplify
$c_{\mathrm{free}}$ locally \\
\addlinespace
Mouse early embryo &
$\sim 10^{3}$ &
progressively consolidating; $\sim$Rouse-like during fate decisions,
TAD-rigid post-fate~\cite{Du2017,Ke2017,Flyamer2017,Saiz2026} &
Energy dissipation / activator proofreading required &
Rescued mainly by long $T$, $k$ redundant enhancers~\cite{Osterwalder2018},
paracrine averaging over $N_{\mathrm{cells}}$~\cite{Etoc2016,Saiz2016}.
CTCF-loop geometry is dispensable for fate specification but required for
post-fate fine-tuning~\cite{Saiz2026,Rao2017,Hsieh2022}. \\
\bottomrule
\end{tabular}
\caption{Two independent axes of fly-vs-mouse asymmetry. The Lammers
axis~\cite{Lammers2023} measures the kinetic burden of distinguishing
cognate from noncognate TFs. The polymer-FPT axis (this work) measures
the structural burden of bringing enhancer and promoter into productive
contact. Each embryo populates the bound~\eqref{eq:star} along both axes
simultaneously.}
\label{tab:two_axes}
\end{table}

The mouse-specific bound on the structural axis becomes, after
multi-enhancer and paracrine rescue,
\begin{equation}
  \left(\frac{\delta c}{c}\right)^{\!2}_{\!\text{mouse, polymer}}
  \;\gtrsim\;
  \frac{\tau_c(\alpha,s_{\min},b)}{k\,T\,N_{\mathrm{cells}}},
  \label{eq:star_mouse}
\end{equation}
where $s_{\min}$ is the effective intra-TAD genomic distance, $k$ is the
number of redundant enhancers contributing in parallel, and
$N_{\mathrm{cells}}$ is the size of the paracrine-averaged cell field. The
mouse minimizes the polymer term by maximizing the prefactors that divide
it; the fly minimizes~$\tau_c$ at the level of the locus itself. Each is a
coherent solution.

The two axes are independent in the sense that an embryo could in principle
sit comfortably on one and at the wire on the other. Empirically, both fly
nc14 and mouse early embryo appear to be at the wire on \emph{both} axes,
which is why morphogen-based developmental decisions are restricted to
specific time windows (early cycles in fly; defined fate-commitment
intervals in mouse) and break down outside them.

\section{Application: 2D mouse gastruloids}
\label{sec:gastruloid}

A 2D mouse gastruloid~\cite{vandenBrink2014,Beccari2018,Moris2020} starts
as a uniform aggregate of pluripotent cells, breaks symmetry under Wnt
activation, and acquires a Brachyury-positive pole opposite a
Sox2-retaining domain. The morphogen gradient is not imposed but
self-organized, through a Wnt/Nodal activator--inhibitor
field~\cite{Muller2012,Etoc2016,Tewary2017}. The total readout-noise
budget acquires an additional source term:
\begin{equation}
  \left(\tfrac{\delta c}{c}\right)^{\!2}_{\!\text{observed}}
  =
  \left(\tfrac{\delta c}{c}\right)^{\!2}_{\!\text{source}}
  +
  \left(\tfrac{\delta c}{c}\right)^{\!2}_{\!\text{propagation}}
  +
  \left(\tfrac{\delta c}{c}\right)^{\!2}_{\!\text{readout}},
  \label{eq:gastruloid_decomp}
\end{equation}
where the readout term is~\eqref{eq:star} populated with mouse parameters.
The gastruloid is a natural test bed for the framework because every term
on the right is, in principle, single-molecule observable: H2B SMT
constrains $\alpha$ and hence $\tau_c$; Sox2 and Bra SMT constrain $D_{\eff}$,
$c_{\mathrm{free}}$, and the kinetic prefactors of the
Lammers~\cite{Lammers2023} sharpness/precision/specificity terms; spatial
statistics of TF bound-fractions across the gastruloid constrain the
source structure.

\subsection{Term-by-term measurement program}

\begin{table}[h]
\centering
\small
\begin{tabular}{p{0.27\linewidth}p{0.30\linewidth}p{0.37\linewidth}}
\toprule
Observable & Term in~\eqref{eq:star} or~\eqref{eq:gastruloid_decomp} &
What it tells you \\
\midrule
H2B SMT ($\alpha$, $\Gamma_\alpha$, fBM state fractions, $P_c$,
return probability) & $\tau_c$ via~\eqref{eq:tau_c} & Whether chromatin is
Rouse-like, crumpled, or TAD-confined; whether $\alpha$ is spatially
patterned across the gastruloid \\
\addlinespace
Sox2 SMT (residence-time distribution; bound fraction; $D_{\eff}$) &
Lammers term (precision; specificity at low $w/c$); $c_{\mathrm{free}}$
input to BP & Pluripotency TF kinetics; baseline against which Bra is
compared \\
\addlinespace
Bra SMT (residence-time distribution; bound fraction; $D_{\eff}$) &
Lammers term at fate transition; $c_{\mathrm{free}}$ input to BP &
Direct gastruloid analog of Bcd readout; expected bimodal residence with
long-residence fraction tracking polarized axis \\
\addlinespace
Spatial autocorrelation of bound fractions across the gastruloid &
Source structure in~\eqref{eq:gastruloid_decomp} & Reaction--diffusion
(intrinsic Turing wavelength) vs.\ localized-source (monotonic, no
wavelength) \\
\addlinespace
Position-dependent variance of Bra bound fraction across replicates &
Relative weight of source vs.\ readout terms in
\eqref{eq:gastruloid_decomp} & Source-dominated: variance position-
independent. Readout-dominated: variance peaks at boundary
(Gregor input/output decomposition) \\
\addlinespace
Time-resolved SMT spanning symmetry breaking & Causal ordering of TF
kinetics and visible polarization & Whether residence-time bimodality
precedes or follows the fate decision \\
\bottomrule
\end{tabular}
\caption{Mapping single-molecule observables to terms in the compound bound
\eqref{eq:star} and the gastruloid noise decomposition
\eqref{eq:gastruloid_decomp}.}
\label{tab:program}
\end{table}

\subsection{Spatial coupling of chromatin and TF readout}

A distinctive feature of gastruloids, absent from fly nc14, is that
chromatin state is itself spatially patterned by the symmetry-breaking
event: Bra-positive cells differentiate and remodel chromatin within hours
of polarization. Hence~$\alpha = \alpha(\mathbf{x})$ across the gastruloid,
$\tau_c = \tau_c(\mathbf{x})$, and the gradient is in two coupled fields ---
a TF concentration field and a chromatin-state field --- rather than one.
Per-cell H2B fBM analysis yields~$\alpha(\mathbf{x})$ directly. A finding of
spatially uniform~$\alpha$ reduces the system to the standard
``TF-reads-concentration'' picture; a finding of patterned~$\alpha$
implicates chromatin in the readout.

\section{Predictions}
\label{sec:predictions}

\begin{enumerate}
  \item \textbf{Chromatin exponent at fly readout loci.} Live two-color
        $eve$ enhancer--promoter MSD in nc14 should yield
        $\alpha = 0.5\pm 0.05$ with no late-time plateau, consistent with
        Br\"uckner~\cite{Bruckner2023}. A value $\alpha\le 0.4$ would require
        the long-range-correlation correction~\cite{Bruckner2023} to lie in
        the favorable direction.

  \item \textbf{Architectural facilitation of pairing at fly loci.} For
        loci with EP separations comparable to or beyond \textit{eve}'s
        142\,kb, mean pairing times in the absence of insulator-mediated
        facilitation should match or exceed Chen \textit{et al.}'s
        $\tau_c\!\sim\!1$\,hr~\cite{Chen2018}. Insulator deletion (e.g.\
        $\Delta$\textit{homie}) should slow pairing further; insulator
        substitution should restore it. Loci with shorter EP separations
        ($\lesssim\!20$\,kb), at which Chen-style architectural
        facilitation is unnecessary, should show pairing times
        approaching the bare-polymer prediction with the
        Br\"uckner~\cite{Bruckner2023} long-range-correlation correction.

  \item \textbf{Transcription-stabilized dwell extension.} Two-color EP
        tracking at fly developmental loci should resolve a bimodal
        distribution of paired-state durations: a short component
        (inactive paired, $\tau_d^{\mathrm{off}}\!\sim\!7$\,min) and a
        long component (active paired, $\tau_d^{\mathrm{on}}\!\sim\!90$\,min)
        consistent with Chen \textit{et al.}~\cite{Chen2018}. Acute
        transcriptional inhibition (e.g.\ triptolide, $\alpha$-amanitin)
        should ablate the long component without much affecting the
        pairing rate, isolating the transcription-coupled dwell extension
        from the architectural pairing mechanism.

  \item \textbf{Mouse loop signature: bimodal contact-duration
        distribution.} Two-color EP tracking at a mouse developmental
        locus (e.g.\ \textit{Sox2}, \textit{Bra}) should show a
        distribution of contact durations that is approximately bimodal:
        a short-dwell component reflecting bare-polymer encounters and a
        long-dwell component reflecting cohesin-anchored loop states with
        $\tau_d \gtrsim \SI{e3}{\second}$~\cite{Gabriele2022}. The
        long-dwell fraction should be small in absolute terms but
        responsible for the bulk of productive transcription. Acute
        cohesin or CTCF depletion~\cite{Saiz2026} should ablate the
        long-dwell component without much affecting the short-dwell rate,
        consistent with the observation that contacts are largely
        preserved~\cite{Hsieh2022} but transcriptional fine-tuning is
        lost~\cite{Saiz2026}.

  \item \textbf{Mouse compensation through enhancer redundancy.} Targeted
        deletion of subsets of redundant enhancers should degrade positional
        precision as~$1/\sqrt{k_{\mathrm{remaining}}}$, where
        $k_{\mathrm{remaining}}$ is the number of functional enhancers
        contributing in parallel.

  \item \textbf{Two-axis decomposition of the fly-vs-mouse split.} The
        Lammers $w/c$ axis~\cite{Lammers2023} and the polymer-FPT $\alpha$
        axis are independent. A hypothetical genome-edited fly with
        artificially elevated $w/c$ but pre-TAD chromatin should fail
        Lammers' specificity bound but pass the polymer-FPT bound; a
        hypothetical fly with low $w/c$ but artificially TAD-consolidated
        chromatin should pass Lammers but fail polymer-FPT. Both should
        fail to pattern correctly, but for different molecular reasons.

  \item \textbf{Gastruloid Bra residence bimodality.} pEM and saSPT
        decomposition of Bra SMT in 2D gastruloids should yield a bimodal
        residence-time distribution with the long-residence fraction
        tracking position along the polarized axis.

  \item \textbf{Gastruloid variance structure.} Cell-to-cell variance of
        Bra free fraction at fixed gastruloid position should be
        position-independent if the source dominates
        \eqref{eq:gastruloid_decomp} and position-dependent (peaked at the
        boundary) if the readout dominates.

  \item \textbf{Gastruloid chromatin patterning.} Per-cell H2B~$\alpha$
        should differ measurably between Bra-high and Sox2-high domains if
        chromatin participates in the readout, and not otherwise.

  \item \textbf{H2B dynamics across acute CTCF/cohesin depletion in
        gastruloids.} The Shi--Shin--Thirumalai HIPPS-DIMES
        framework~\cite{Shi2025} predicts that loss of cohesin reorganizes
        the chromatin connectivity matrix~$\mathbf{K}$ --- redistributing
        locus centralities $C_i$ and closing the low-mode gap --- while
        leaving the local fractional-Langevin exponent~$\alpha$ largely
        intact. The expected single-molecule signatures during the
        48--72\,h fate-decision window are therefore (i) a modest
        ensemble-mean increase ($\sim$20--40\%) in H2B effective
        diffusivity~$D$ with no detectable shift in~$\alpha$; (ii) a
        locus-by-locus $\Delta D$ that anticorrelates with the relative
        change in centrality $C_i^{\Delta\mathrm{RAD21}}/C_i^{\mathrm{WT}}$
        computed from matched Micro-C; (iii) a redistribution of pEM
        state fractions reflecting the centrality reorganization, with
        within-state~$\alpha$ preserved. This prediction is consistent
        with Mazzocca \textit{et al.}~\cite{Mazzocca2024} MINFLUX+HILO
        measurements (small effect on~$D$, no shift in~$\alpha$ on RAD21
        depletion) and with pEM analyses of H2B in 3617 cells under RAD21
        depletion that show shifted state fractions without a within-state
        $\alpha$ shift. At 120--168\,h, when the looping function becomes
        essential~\cite{Saiz2026}, locus mobility and pEM state weights
        should diverge between WT and $\Delta$N-term CTCF gastruloids by a
        larger margin. A finding that within-state~$\alpha$ drops sharply
        at 48--72\,h under cohesin or CTCF loss would contradict both
        Mazzocca \textit{et al.}~\cite{Mazzocca2024} and the HIPPS-DIMES
        prediction, and force a revision of the mouse-side picture.

  \item \textbf{Signed locus-specific~$\tau$ change inside vs.\ at TAD
        boundaries.} Shi \textit{et al.}~\cite{Shi2025} predict that
        cohesin depletion changes the two-point relaxation
        time~$\tau_{ij}$ in opposite directions for intra-TAD pairs (slower,
        because $\langle r_{ij}\rangle$ grows) and TAD-boundary pairs
        (faster). Two-color EP tracking in CTCF-AID or RAD21-AID
        gastruloids should resolve this signed, locus-specific shift. The
        prediction is consistent with the Saiz \textit{et al.}
        morphogenesis-vs-fate dissociation~\cite{Saiz2026}: morphogenetic
        precision degrades because $\tau_c$ is differentially detuned
        across loci, not because $\tau_c$ uniformly worsens, while
        time-averaged contact probabilities (and therefore
        threshold-crossing fate decisions) remain robust.

  \item \textbf{Bra/Sox2 residence times across CTCF depletion.} Bra and
        Sox2 SMT in CTCF-AID gastruloids should show comparable residence-
        time distributions and bound fractions at 48--72\,h with and
        without CTCF, given that fate specification is preserved. Any
        appreciable difference would localize the early CTCF requirement
        to TF binding kinetics rather than to specific gene
        regulation~\cite{Saiz2026}.

  \item \textbf{Substrate stiffness as a knob on $\tau_c$.} Stiffer
        substrates should lower~$\alpha(\mathrm{H2B})$, raise~$\tau_c$, and
        degrade boundary precision in proportion to~\eqref{eq:star_mouse}.
        Boundary precision unchanged by stiffness implies the readout
        operates with margin on the polymer axis (the Lammers axis may still
        be tight).
\end{enumerate}

\section{Discussion}

The Gregor 2007 result was that \textit{Drosophila} development reads
positional information at the Berg--Purcell limit. Lammers \textit{et al.}
2023 added that the same readout is bounded by a kinetic precision/
sharpness/specificity tradeoff, with fly and mouse on opposite sides of the
$w/c$ axis. The polymer-revised statement here adds a third constraint of
comparable weight: chromatin first-passage time. The three constraints are
physically independent --- diffusive counting (Berg--Purcell), gene-circuit
kinetics under noncognate interference (Lammers), and polymer first-passage
(this work) --- and a developmental readout must satisfy all three
simultaneously. The Xiao--Hafner--Boettiger~\cite{Xiao2021} futile-cycle
amplifier sits downstream as a high-gain decoder that converts small
systematic differences (in $c$ or in $\tau_c$) into sharp transcriptional
responses, while equally amplifying any residual stochastic noise; the
biology lives in the regime where signal-amplification mechanisms operate
on inputs whose stochastic component has been independently suppressed.

\textit{Drosophila} nc14 satisfies all three because $w/c$ is low,
chromatin is pre-TAD with $\alpha\!\approx\!0.5$, and the syncytium permits
spatial averaging. Mouse early embryos satisfy all three through entirely
different mechanisms: energy dissipation rescues the Lammers term;
multi-enhancer redundancy and paracrine averaging rescue the polymer term;
hours-scale fate-commitment integration windows accommodate the residual
Berg--Purcell term; and the futile-cycle gain~\cite{Xiao2021} converts the
small intra-TAD vs.\ inter-TAD contact-frequency differences into the
$\sim$10-fold expression boundaries actually observed. The 2D gastruloid
offers the cleanest experimental setting in which to measure each term
independently, because the geometry is finite, the chromatin is mammalian,
the symmetry-breaking event is controllable, and every term
in~\eqref{eq:star} except $\tau_c$ itself is single-molecule observable.

\paragraph{Caveats.} The bound~\eqref{eq:star} assumes the renewal
approximation for EP contacts and weak correlation between TF occupancy and
contact events. The polymer scaling~\eqref{eq:tau_c} assumes equilibrium
fractional Langevin dynamics; non-equilibrium driving (cohesin loop
extrusion, RNA Pol II--mediated stiffening) violates~\eqref{eq:tau_c} and
is plausibly the origin of the Br\"uckner~\cite{Bruckner2023} long-range
correlations. We do \emph{not} invoke condensate-mediated capture as a
mechanism: while Sox2 and other developmental TFs form transient
clusters, direct evidence that condensation reduces $\tau_c$ in a
fate-relevant way is lacking. The role of structural loops in setting
$\tau_c$ is also weaker than the rigid-TAD interpretation suggests: acute
cohesin and CTCF
perturbations~\cite{Rao2017,Schwarzer2017,Hsieh2022,Saiz2026} preserve
substantial EP contact frequency and, in gastruloids, preserve fate
specification entirely. Shi \textit{et al.}~\cite{Shi2025} sharpen this
picture: the rapid, locus-specific dynamics that distinguishes WT
chromatin from a polymer-shuffled control is encoded in the
\emph{full} static contact map (the connectivity matrix~$\mathbf{K}$
derivable from Micro-C), not in CTCF loops alone. The framework is
therefore most useful when read as constraining the static
contact-map-derived dynamics --- captured empirically by per-cell H2B SMT
diffusivity together with matched Micro-C-derived locus centralities,
and consistent with the Mazzocca \textit{et al.}~\cite{Mazzocca2024}
finding that within-state~$\alpha$ is preserved under cohesin loss while
$D$ shifts modestly --- rather than as constraining the engineered loop
topology imposed by CTCF/cohesin specifically. The formal coupling between the
Lammers term and the polymer term --- in particular, whether locus
localization to general transcription machinery should enter as a separate
$\tau_c$ or as an additional Lammers activation step~\cite{Lammers2023} ---
is left for future work.

\begin{thebibliography}{99}

\bibitem{Lammers2023}
N.~C.~Lammers, A.~I.~Flamholz, H.~G.~Garcia. ``Competing constraints shape
the nonequilibrium limits of cellular decision-making.''
\textit{PNAS} \textbf{120}, e2211203120 (2023).

\bibitem{Bruckner2023}
D.~B.~Br\"uckner, H.~Chen, L.~Barinov, B.~Zoller, T.~Gregor. ``Stochastic
motion and transcriptional dynamics of pairs of distal DNA loci on a
compacted chromosome.'' \textit{Science} \textbf{380}, 1357--1362 (2023).

\bibitem{Shi2025}
G.~Shi, S.~Shin, D.~Thirumalai. ``Static three-dimensional structures
determine fast dynamics between distal loci pairs in interphase
chromosomes.'' \textit{bioRxiv} 2025.01.16.633484 (2025).

\bibitem{Garcia2021NAR}
D.~A.~Garcia, G.~Fettweis, D.~M.~Presman, V.~Paakinaho, C.~Jarzynski,
A.~Upadhyaya, G.~L.~Hager. ``Power-law behavior of transcription factor
dynamics at the single-molecule level implies a continuum of affinities.''
\textit{Nucleic Acids Research} \textbf{49}, 6605--6620 (2021).

\bibitem{KimLarson2025}
J.~M.~Kim, D.~R.~Larson. ``Timing is everything: transcription bursting
in development.'' \textit{Genes \& Development} \textbf{39}, 995--1011
(2025).

\bibitem{Mazzocca2024}
M.~Mazzocca \textit{et al.} ``Live-cell MINFLUX and HILO imaging of
histone H2B dynamics under acute cohesin depletion.'' (Hansen lab,
\textit{ref.\ to be completed}).

\bibitem{Lucas2018}
T.~Lucas, H.~Tran, C.~A.~Perez Romero, A.~Guillou, C.~Fradin, M.~Coppey,
A.~M.~Walczak, N.~Dostatni. ``3 minutes to precisely measure morphogen
concentration.'' \textit{PLoS Genetics} \textbf{14}, e1007676 (2018).

\bibitem{Chen2018}
H.~Chen, M.~Levo, L.~Barinov, M.~Fujioka, J.~B.~Jaynes, T.~Gregor.
``Dynamic interplay between enhancer--promoter topology and gene
activity.'' \textit{Nat.\ Genet.} \textbf{50}, 1296--1303 (2018).

\bibitem{Bothma2014}
J.~P.~Bothma \textit{et al.} ``Dynamic regulation of \textit{eve} stripe~2
expression reveals transcriptional bursts in living \textit{Drosophila}
embryos.'' \textit{PNAS} \textbf{111}, 10598--10603 (2014).

\bibitem{Lammers2020}
N.~C.~Lammers \textit{et al.} ``Multimodal transcriptional control of
pattern formation in embryonic development.''
\textit{PNAS} \textbf{117}, 836--847 (2020).

\bibitem{Saiz2026}
N.~Alonso Saiz, M.~Martinovic, M.~Rubio, P.~Samal, S.~Giselbrecht,
L.~Braccioli, E.~de Wit. ``A dual role for CTCF in development.''
\textit{bioRxiv} 2026.03.17.712290 (2026).

\bibitem{Rao2017}
S.~S.~P.~Rao \textit{et al.} ``Cohesin loss eliminates all loop domains.''
\textit{Cell} \textbf{171}, 305--320 (2017).

\bibitem{Schwarzer2017}
W.~Schwarzer \textit{et al.} ``Two independent modes of chromatin
organization revealed by cohesin removal.''
\textit{Nature} \textbf{551}, 51--56 (2017).

\bibitem{Hsieh2022}
T.-H.~S.~Hsieh \textit{et al.} ``Enhancer--promoter interactions and
transcription are largely maintained upon acute loss of CTCF, cohesin,
WAPL or YY1.'' \textit{Nat.\ Genet.} \textbf{54}, 1919--1932 (2022).

\bibitem{Xiao2021}
J.~Y.~Xiao, A.~Hafner, A.~N.~Boettiger. ``How subtle changes in 3D
structure can create large changes in transcription.''
\textit{eLife} \textbf{10}, e64320 (2021).

\bibitem{Kaizu2014}
K.~Kaizu, W.~de Ronde, J.~Paijmans, K.~Takahashi, F.~Tostevin,
P.~R.~ten Wolde. ``The Berg--Purcell limit revisited.''
\textit{Biophysical J.} \textbf{106}, 976--985 (2014).

\bibitem{Gregor2007a}
T.~Gregor, D.~W.~Tank, E.~F.~Wieschaus, W.~Bialek. ``Probing the limits to
positional information.'' \textit{Cell} \textbf{130}, 153--164 (2007).

\bibitem{Gregor2007b}
T.~Gregor, E.~F.~Wieschaus, A.~P.~McGregor, W.~Bialek, D.~W.~Tank.
``Stability and nuclear dynamics of the Bicoid morphogen gradient.''
\textit{Cell} \textbf{130}, 141--152 (2007).

\bibitem{Bialek2005}
W.~Bialek, S.~Setayeshgar. ``Physical limits to biochemical signaling.''
\textit{PNAS} \textbf{102}, 10040--10045 (2005).

\bibitem{Amitai2013}
A.~Amitai, D.~Holcman. ``Polymer model with long-range interactions.''
\textit{Phys.\ Rev.\ E} \textbf{88}, 052604 (2013).

\bibitem{Guerin2016}
T.~Gu\'erin, O.~B\'enichou, R.~Voituriez. ``Mean first-passage times of
non-Markovian random walkers in confinement.'' \textit{Nature} \textbf{534},
356--359 (2016).

\bibitem{Gabriele2022}
M.~Gabriele \textit{et al.} ``Dynamics of CTCF- and cohesin-mediated
chromatin looping revealed by live-cell imaging.''
\textit{Science} \textbf{376}, 496--501 (2022).

\bibitem{Mach2022}
P.~Mach \textit{et al.} ``Cohesin and CTCF control the dynamics of
chromosome folding.'' \textit{Nat.\ Genet.} \textbf{54}, 1907--1918 (2022).

\bibitem{Mir2017}
M.~Mir \textit{et al.} ``Dense Bicoid hubs accentuate binding along the
morphogen gradient.'' \textit{Genes \& Dev.} \textbf{31}, 1784--1794 (2017).

\bibitem{Mir2018}
M.~Mir \textit{et al.} ``Dynamic multifactor hubs interact transiently with
sites of active transcription in \textit{Drosophila} embryos.''
\textit{eLife} \textbf{7}, e40497 (2018).

\bibitem{Hug2017}
C.~B.~Hug, A.~G.~Grimaldi, K.~Kruse, J.~M.~Vaquerizas. ``Chromatin
architecture emerges during zygotic genome activation independent of
transcription.'' \textit{Cell} \textbf{169}, 216--228 (2017).

\bibitem{Ogiyama2018}
Y.~Ogiyama, B.~Schuettengruber, G.~L.~Papadopoulos, J.-M.~Chang,
G.~Cavalli. ``Polycomb-dependent chromatin looping contributes to gene
silencing during \textit{Drosophila} development.''
\textit{Mol.\ Cell} \textbf{71}, 73--88 (2018).

\bibitem{Du2017}
Z.~Du \textit{et al.} ``Allelic reprogramming of 3D chromatin architecture
during early mammalian development.'' \textit{Nature} \textbf{547},
232--235 (2017).

\bibitem{Ke2017}
Y.~Ke \textit{et al.} ``3D chromatin structures of mature gametes and
structural reprogramming during mammalian embryogenesis.''
\textit{Cell} \textbf{170}, 367--381 (2017).

\bibitem{Flyamer2017}
I.~M.~Flyamer \textit{et al.} ``Single-nucleus Hi-C reveals unique chromatin
reorganization at oocyte-to-zygote transition.''
\textit{Nature} \textbf{544}, 110--114 (2017).

\bibitem{Osterwalder2018}
M.~Osterwalder \textit{et al.} ``Enhancer redundancy provides phenotypic
robustness in mammalian development.''
\textit{Nature} \textbf{554}, 239--243 (2018).

\bibitem{Saiz2016}
N.~Saiz \textit{et al.} ``Asynchronous fate decisions by single cells
collectively ensure consistent lineage composition in the mouse blastocyst.''
\textit{Nat.\ Commun.} \textbf{7}, 13463 (2016).

\bibitem{Etoc2016}
F.~Etoc \textit{et al.} ``A balance between secreted inhibitors and edge
sensing controls gastruloid self-organization.''
\textit{Dev.\ Cell} \textbf{39}, 302--315 (2016).

\bibitem{Muller2012}
P.~M\"uller \textit{et al.} ``Differential diffusivity of Nodal and Lefty
underlies a reaction-diffusion patterning system.''
\textit{Science} \textbf{336}, 721--724 (2012).

\bibitem{Tewary2017}
M.~Tewary \textit{et al.} ``A stepwise model of reaction-diffusion and
positional information governs self-organized human peri-gastrulation-like
patterning.'' \textit{Development} \textbf{144}, 4298--4312 (2017).

\bibitem{vandenBrink2014}
S.~C.~van den Brink \textit{et al.} ``Symmetry breaking, germ layer
specification and axial organisation in aggregates of mouse embryonic stem
cells.'' \textit{Development} \textbf{141}, 4231--4242 (2014).

\bibitem{Beccari2018}
L.~Beccari \textit{et al.} ``Multi-axial self-organization properties of
mouse embryonic stem cells into gastruloids.''
\textit{Nature} \textbf{562}, 272--276 (2018).

\bibitem{Moris2020}
N.~Moris \textit{et al.} ``An in vitro model of early anteroposterior
organization during human development.''
\textit{Nature} \textbf{582}, 410--415 (2020).

\end{thebibliography}

\end{document}
